\documentclass[12pt]{article}
\usepackage[margin=1in]{geometry}
\usepackage{amsmath,amssymb,amsthm}
\usepackage{fontspec}
\setmainfont{Noto Serif CJK JP}
\XeTeXlinebreaklocale "ja"
\XeTeXlinebreakskip = 0pt plus 1pt minus 0.1pt
\usepackage{hyperref}
\usepackage{enumitem}
\usepackage{booktabs}
\usepackage{caption}

\theoremstyle{plain}
\newtheorem{theorem}{定理}[section]
\newtheorem{proposition}[theorem]{命題}
\newtheorem{corollary}[theorem]{系}
\theoremstyle{definition}
\newtheorem{definition}[theorem]{定義}
\newtheorem{remark}[theorem]{注意}
\numberwithin{equation}{section}

\newcommand{\R}{\mathbb{R}}
\newcommand{\C}{\mathbb{C}}
\newcommand{\Z}{\mathbb{Z}}
\newcommand{\HH}{\mathbb{H}}
\newcommand{\br}{\mathrm{br}}
\newcommand{\vis}{\mathrm{vis}}
\newcommand{\eff}{\mathrm{eff}}
\newcommand{\LoN}{\mathrm{LoN}}
\newcommand{\Pl}{\mathrm{Pl}}
\newcommand{\fU}{\mathfrak{U}}
\newcommand{\fI}{\mathfrak{I}}
\newcommand{\fV}{\mathfrak{V}}
\newcommand{\fB}{\mathfrak{B}}
\DeclareMathOperator{\Tr}{Tr}
\DeclareMathOperator{\Gr}{Gr}
\DeclareMathOperator{\sech}{sech}
\renewcommand{\proofname}{証明}

\begin{document}

\begin{titlepage}
\centering
\vspace*{1cm}
{\LARGE\bfseries 余剰次元の脱リテラル化：\\[6pt]
コンパクト化された時空座標の\\[6pt]
4D代替としての$X(2)$上の内部モジュラー幾何}
\vspace{1.5cm}

{\large 中村 昌義$^{1,*}$}
\vspace{0.8cm}

{\normalsize
$^1$\,北与野メンタルクリニック、〒338-0002
埼玉県さいたま市中央区下落合4-20-14\\
\texttt{mjolnir501@gmail.com}\\[4pt]
$^*$\,責任著者
}
\vspace{0.8cm}

{\small
\noindent\textbf{関連リポジトリ：}\\
LoNalogyフレームワーク (v87.1):
\href{https://doi.org/10.5281/zenodo.19115338}{doi:10.5281/zenodo.19115338}\\
Yang--Mills証明 (v90):
\href{https://doi.org/10.5281/zenodo.19183838}{doi:10.5281/zenodo.19183838}\\
姉妹論文：$X(2)$上の統計力学、素粒子宇宙論 (v9)、
ハッブル再構成 (v2)、コロナ加熱
}
\vspace{1.0cm}

\begin{abstract}
\noindent
弦理論は量子重力の整合性のために余剰空間次元を要求する。
LHC（ADD重力子放出に対して
$M_D\gtrsim 5.5$--$11.2\,\mathrm{TeV}$、
RSカルツァ-クライン重力子に対して
$m_1\gtrsim 2.2$--$5.6\,\mathrm{TeV}$）
および短距離重力実験（$\delta=1$で$R<44\,\mu\mathrm{m}$）
による広範な探索にもかかわらず、
余剰次元の証拠は見つかっていない。

LoNalogyフレームワークではリテラルな余剰時空座標が不要であることを示す。
核心作用は4次元であり、弦理論で余剰次元が果たす役割——ゲージ統一、
世代構造、結合階層——は$X(2)$の内部モジュラー/ファイバー幾何が担う。
LoNalogyに現れる整数$\nu_{\br}=12$、$d_{\mathrm{fib}}=7$、$N=5$は
内部束の次元であり時空座標ではない。

結論：余剰次元は否定されるのではなく\emph{脱リテラル化}される。
4D法則の背後にある内部幾何として再解釈され、
これはワープ型余剰次元を4D複合ダイナミクスとして読む
主流のAdS/CFT的解釈と整合する。
\end{abstract}
\end{titlepage}

\setcounter{page}{1}
\tableofcontents
\newpage


\section{序論}\label{sec:intro}

カルツァ（1921年）とクライン（1926年）以来、時空が4次元以上を
持つ可能性は理論物理学で繰り返し主題となってきた。
弦理論はアノマリー消去と量子整合性のために10次元または11次元の
時空を要求する。余剰次元は通常スケール$R$でコンパクト化され、
カルツァ-クライン（KK）励起が質量$m_n\sim n/R$で現れると仮定される。

実験的に研究された三つの広いクラスがある：
(i) ADD（大きな平坦余剰次元、基本スケール$M_D$がTeVスケール）、
(ii) RS（一つのワープ余剰次元、指数的階層）、
(iii) UED（標準模型場がバルクに伝播）。

広範な探索にもかかわらず、いずれの証拠も見つかっていない。


\section{実験的ステータス}\label{sec:experiment}

\subsection{LHCの制約}

LHC Run~2データセット（$\sqrt{s}=13\,\mathrm{TeV}$、
$\sim140\,\mathrm{fb}^{-1}$）は余剰次元モデルに対する
最強の加速器制約を与える。

\begin{center}
\begin{tabular}{lll}
\toprule
モデル & チャンネル & 制約\\
\midrule
ADD（重力子放出）& モノジェット & $M_D\gtrsim 5.5$--$11.2\,\mathrm{TeV}$\\
ADD（仮想重力子）& ダイジェット角度 & $M_{TT}>9.02\,\mathrm{TeV}$\\
RS（KK重力子）& ダイフォトン & $m_1\gtrsim 2.2$--$5.6\,\mathrm{TeV}$\\
ラディオン & ATLAS/CMS & $m_{\mathrm{radion}}\gtrsim 3.1$--$3.2\,\mathrm{TeV}$\\
微小BH & 高多重度 & $M_{\mathrm{BH}}\lesssim 10.1\,\mathrm{TeV}$排除\\
UED & 各種 & $1/R\gtrsim 1.4$--$1.5\,\mathrm{TeV}$\\
\bottomrule
\end{tabular}
\end{center}

最小UEDはLHC制約と暗黒物質レリック密度の要件を合わせると排除される。

\subsection{短距離重力実験}

卓上重力実験は大きな平坦余剰次元を直接制約する。
$\delta=1$の平坦余剰次元に対して$R<44\,\mu\mathrm{m}$、
$\delta=2$ではATLASモノジェットから$R<3.8\,\mu\mathrm{m}$。

\subsection{実験的ステータスの要約}

\begin{proposition}[実験的判定]\label{prop:experiment}
\begin{equation}
\boxed{\text{低スケールのリテラルな余剰次元：排除方向に強く制約。}}
\end{equation}
\begin{equation}
\boxed{\parbox{0.8\textwidth}{\centering
プランクスケールでコンパクト化された余剰次元：\\
論理的には可能だが実験的に到達不可能。}}
\end{equation}
\end{proposition}


\section{LoNalogyは4D理論である}\label{sec:4D}

\subsection{核心作用}

\begin{definition}[LoNalogy核心作用]\label{def:action}
LoNalogyのマスター作用は4次元時空上で定義される：
\begin{equation}\label{eq:action}
S_{\LoN} = S_{\mathrm{grav}} + S_x + S_{\vis} + S_A + S_X^{\eff},
\end{equation}
ここで
\begin{equation}
S_{\mathrm{grav}}+S_x
= \int d^4x\,\sqrt{-g}\left[
\frac{M_{\Pl}^2}{2}R
-\frac{1}{2}G_{xx}(x)\,\partial_\mu x\,\partial^\mu x
-V(x)\right].
\end{equation}
可視、ポータル、ブリッジの各セクターも同様に4D作用として書かれる。
第5以上の座標$x^5,x^6,\ldots$は現れない。
\end{definition}

\subsection{内部次元 vs 時空次元}

LoNalogyフレームワークにはいくつかの大きな整数が現れる。
これらを時空次元と区別することが不可欠である。

\begin{proposition}[内部次元は時空座標ではない]\label{prop:internal}
LoNalogyに現れる整数の起源：

(i) $\nu_{\br}=12$：破れた接空間
$SU(5)/S(U(3)\times U(2))\cong\Gr_2(\C^5)$の実次元
（$2\cdot 3\cdot 2=12$）。内部コセット次元であり時空座標数ではない。

(ii) $d_{\mathrm{fib}}=7$：三セクター分解のファイバー/ブリッジ次元。
内部自由度の数であり空間方向ではない。

(iii) $N=5$（ランク）：湯川閉包
$\Lambda^2 V\otimes\Lambda^2 V\otimes V\to\det V$と
アノマリー消去から選択。束のランクであり時空次元ではない。
\end{proposition}

\begin{remark}
内部次元と時空次元の混同はLoNalogyの最も一般的な誤読である。
弦理論に馴染みのある者は、内部次元が\emph{リテラルに}
余分な時空座標（コンパクト化された）であることに慣れている。
LoNalogyではそうではない。
\end{remark}


\section{余剰次元の脱リテラル化}\label{sec:delit}

\begin{theorem}[脱リテラル化]\label{thm:delit}
弦/KK理論で余剰次元が果たす全ての構造的役割は、
LoNalogyに4D対応物を持つ：

\begin{center}
\begin{tabular}{lll}
\toprule
役割 & 弦/KK言語 & LoNalogy言語\\
\midrule
ゲージ統一 & $G/H$上のコンパクト化 & $X(2)$上のランク5束\\
世代構造 & 固定点/交差 & 3カスプ$\{0,1,\infty\}$\\
結合階層 & ワープ因子$e^{-kr_c\pi}$ & ブリッジ相互律\\
KKタワー & $m_n = n/R$ & ホロノミーフーリエモード\\
モジュライ安定化 & フラックスコンパクト化 & strict kernel方程式\\
隠れたセクター & バルク/ブレーン分離 & $\fI$-セクター、Schur補行列\\
\bottomrule
\end{tabular}
\end{center}

各ケースでLoNalogyの構成は完全に4D内で、
内部（ファイバー/モジュラー）構造とともに生きており、
追加の時空座標を導入しない。
\end{theorem}

\begin{proof}
各行はLoNalogyフレームワークとの直接比較により検証される。

\emph{ゲージ統一。}
弦理論はアイソメトリー群がゲージ群を与える多様体$G/H$上で
コンパクト化する。LoNalogyはアノマリーフリーな分解$5=3+2$を通じて
$X(2)$上のランク5束から$SU(3)\times SU(2)\times U(1)$を得る。

\emph{世代構造。}
弦理論では3世代はコンパクト多様体の交差数や固定点から生じる。
LoNalogyでは3世代は$X(2)$の3カスプ$\{0,1,\infty\}$に対応する。

\emph{結合階層。}
RSモデルはワープ因子$e^{-kr_c\pi}$でプランク/TeV階層を生成する。
LoNalogyはブリッジ相互律$M_C\|\sigma_{\br}\|=\Lambda_*$と
strict kernel方程式で階層を代数的に固定する。

\emph{KKタワー。}
KK理論ではタワー$m_n=n/R$はコンパクト方向のフーリエ分解から生じる。
LoNalogyではホロノミー変数$u\in[0,1)$が類似の役割を果たすが、
時空座標ではなくファイバー座標である。

\emph{モジュライ安定化。}
弦理論はフラックスコンパクト化でモジュライを固定する。
LoNalogyはstrict kernel方程式$\nu_{\br}(1-k)^2(1+k)=2$で
モジュライを固定する。

\emph{隠れたセクター。}
ブレーンワールドモデルでは隠れたセクターは余剰次元で
分離された別のブレーン上に住む。LoNalogyでは隠れたセクター$\fI$は
空間的距離ではなくSchur補行列構造によって分離される。
\end{proof}

\begin{corollary}[余剰次元の脱リテラル化]
\begin{equation}
\boxed{\text{余剰次元は否定されるのではなく脱リテラル化される。}}
\end{equation}
余剰次元の構造的内容——統一、世代構造、階層、タワー、
モジュライ、隠れたセクター——はLoNalogyで保存されるが、
追加の時空座標ではなく内部モジュラー/ファイバー幾何に収容される。
\end{corollary}


\section{AdS/CFTとの整合性}\label{sec:adscft}

\begin{proposition}[ホログラフィック双対性との両立]\label{prop:adscft}
LoNalogyにおける余剰次元の脱リテラル化はAdS/CFT対応と整合し、
実際にそれによって予見されている。
RSモデルでは5Dワープ幾何$\mathrm{AdS}_5$は
AdS/CFT辞書を通じて4D強結合複合理論と双対である。
LoNalogyはこの双対性の「4D側」を最初から選択する。
\end{proposition}


\section{LHCが探すべきもの}\label{sec:lhc}

\begin{proposition}[LoNalogy予測 vs KKタワー]\label{prop:lhc}
LoNalogyが正しければ、LHCはKKタワー（$m_n=n/R$）を
見ることを期待すべきではない。新物理のシグナチャーは
4Dブリッジ/ポータル/隠れたセクター現象である：

(i) 暗黒物質候補$\chi$（$338$--$348\,\mathrm{GeV}$、ブリッジセクター）。

(ii) 軽いポータルスカラー$A'$
（$m_{A'}=0.204\,\mathrm{GeV}$）。

(iii) LoNalogy予測レートでの陽子崩壊。

これらは4D予測であり余剰次元のシグナチャーではない。
\begin{equation}
\boxed{\text{正しい問い：LHCは4Dブリッジ/ポータル構造を見られるか？}}
\end{equation}
\end{proposition}


\section{最終判定}\label{sec:verdict}

\begin{center}
\begin{tabular}{p{5.5cm}l}
\toprule
命題 & ステータス\\
\midrule
リテラルな余剰時空次元が実在する & 未確認\\
低スケール（$\sim$TeV）余剰次元 & 強く制約\\
プランクスケールのコンパクト化 & 論理的に可能、実験的に到達不可能\\
LoNalogyがリテラルな余剰次元を要求する & \textbf{いいえ}\\
LoNalogyが余剰次元の構造的役割を再現する & \textbf{はい}（内部幾何を通じて）\\
\bottomrule
\end{tabular}
\end{center}

最短の答え：
\begin{equation}
\boxed{\text{余剰次元：否定ではなく脱リテラル化。LoNalogyは4D完備。}}
\end{equation}


\section*{結論}
\addcontentsline{toc}{section}{結論}

実験的ステータスは明白：余剰次元の証拠は見つかっておらず、
最も単純な低スケールモデルは次第に排除されつつある。
LoNalogyフレームワークはこのヌル結果の説明を与える：
リテラルな余剰時空座標は最初から不要だった。

余剰次元の構造的内容——ゲージ統一、世代構造、結合階層、
モジュライ安定化、隠れたセクター——は$X(2)$の内部
モジュラー/ファイバー幾何に完全に捕捉される。
核心作用は4次元である。

一文で：

\begin{quote}
\emph{余剰次元とは、内部モジュラー幾何が時空座標の言語の上に
投じる影である。}
\end{quote}


\section*{謝辞}
\addcontentsline{toc}{section}{謝辞}

著者はClaude (Anthropic)に、LoNalogyフレームワークと
余剰次元モデルとの関係の共同分析を感謝する。


\begin{thebibliography}{99}

\bibitem{LoNalogyv87}
M.~Nakamura, \textit{LoNalogy v87.1},
Zenodo (2026),
\href{https://doi.org/10.5281/zenodo.19115338}{doi:10.5281/zenodo.19115338}.

\bibitem{YMmassGap}
M.~Nakamura, \textit{Yang--Mills Existence and Mass Gap},
Zenodo (2026),
\href{https://doi.org/10.5281/zenodo.19183838}{doi:10.5281/zenodo.19183838}.

\bibitem{PDG2025}
R.~L.~Workman \textit{et al.}\ (Particle Data Group),
\textit{Review of Particle Physics},
Prog.\ Theor.\ Exp.\ Phys.\ (2025).

\bibitem{ADD}
N.~Arkani-Hamed, S.~Dimopoulos, and G.~Dvali,
Phys.\ Lett.\ B \textbf{429}, 263 (1998).

\bibitem{RS}
L.~Randall and R.~Sundrum,
Phys.\ Rev.\ Lett.\ \textbf{83}, 3370 (1999).

\end{thebibliography}

\end{document}
